% ============================================================
% CHAPTER 4
% ============================================================

\chapter{L-Systems as Sigma Models}
\label{chap:lsystems-sigmamodels}

\section{Motivation}

Lindenmayer systems (L--systems) provide a generative formalism for
modeling biological growth and recursive geometric structures. In the
context of the lamphrodynamic plenum, L--systems become dynamical rules
for rewriting geometric configurations. The key idea of this chapter is
that L--systems admit a natural interpretation as \emph{sigma models}:
their rules are maps into a target space of geometric patterns, and
their rewriting dynamics evolve as field configurations on a worldsheet.

This viewpoint enables a derived geometrical interpretation of
L--systems, with rewrite rules realized as points of a derived moduli
stack $R$, and rewriting steps realized as lamphrodynamic evolutions.
Quantization via the AKSZ/BV formalism then yields a quantum theory of
semantic evolution.

\section{L-Systems}

An (untyped) L--system consists of an alphabet $\Sigma$, a set of
production rules $P\colon \Sigma\to \Sigma^\ast$, and an initial word
$\omega\in\Sigma^\ast$. The rewriting dynamics generate a sequence
$\omega, P(\omega), P(P(\omega)),\dots$, whose geometric interpretation
depends on assigning geometric primitives to symbols in $\Sigma$.

\begin{definition}
An \emph{entropic L--system} is an L--system whose alphabet $\Sigma$
is equipped with lamphrodynamic weights and boundary conditions, and for
which production rules preserve entropy admissibility.
\end{definition}

\begin{remark}
Entropic L--systems represent lamphrodynamics constrained by recursive
semantic structure. They provide a bridge between biological growth,
computational generation, and lamphrodynamic dynamics.
\end{remark}

\section{Sigma Model Interpretation}

Let $X$ be a two--dimensional worldsheet representing the space of
rewriting steps. Let $\mathcal{T}$ denote the derived target stack of
geometric primitives. Then an L--system can be interpreted as a map
\[
  \phi \colon X \longrightarrow \mathcal{T},
\]
where the rewriting rule supplies boundary conditions on $\partial X$.

\begin{definition}
The \emph{L--system sigma model} is the field theory whose fields are
maps $\phi\colon X\to\mathcal{T}$ determined by the production rules of
the L--system, and whose action functional encodes lamphrodynamic
entropy.
\end{definition}

\section{Derived Moduli of Rules}

Each production rule $a\mapsto w\in\Sigma^\ast$ determines a geometric
constraint on fields $\phi$ at the relevant boundary components. The
collection of such rules forms a derived moduli stack
\[
  R := \{\text{admissible production rules}\}.
\]

\begin{lemma}
$R$ is a derived Artin stack governed by lamphrodynamic boundary
conditions and derived gluing rules, with tangent complex determined by
variations of boundary entropy.
\end{lemma}

\begin{proof}[Sketch]
Each rule $a\mapsto w$ specifies a derived matching condition between
boundary components of the configuration. Derived gluing yields an Artin
stack, and lamphrodynamic admissibility ensures that tangent complexes
exist with appropriate degeneration.
\end{proof}

\section{BV Quantization}

To quantize the L--system sigma model, we apply the AKSZ/BV formalism
(see Volume~I, Chapters~4 and~7). The source is the worldsheet $X$, and
the target is the derived stack $R$ of production rules.

\begin{definition}
The \emph{BV--quantized L--system} is the AKSZ sigma model with
source $X$ and target $R$, whose BV action is constructed from
lamphrodynamic constraints on $\phi$.
\end{definition}

\begin{proposition}
The BV action satisfies the classical master equation
$\{S,S\}=0$, and defines a BV complex whose cohomology classifies
lamphrodynamic deformations of admissible rewrite rules.
\end{proposition}

\begin{proof}[Sketch]
AKSZ guarantees the master equation whenever the target carries a
$(-1)$--shifted symplectic structure. The derived stack $R$ inherits
such a structure from the lamphrodynamic entropy form on $\mathcal{T}$
by transgression.
\end{proof}

\section{Worked Example: Algae L-System}

Consider the classical algae L--system with alphabet $\{A,B\}$ and
production rules
\[
  A\mapsto AB,\qquad B\mapsto A.
\]
Assign entropy weights $s_A,s_B$ to the symbols and interpret $A,B$ as
basic geometric primitives subject to lamphrodynamic boundary
conditions.

\begin{proposition}
The BV--quantized algae system realizes a lamphrodynamic sigma model
whose field content corresponds to alternating lamphrodynamic boundary
entropies $s_A$ and $s_B$, and whose BV action reduces to a sum of
boundary contributions.
\end{proposition}

\begin{proof}[Sketch]
Each production $A\mapsto AB$ attaches boundary conditions corresponding
to $A$ and $B$ in sequence. Collapse in the BV complex reduces internal
structure to boundary alternation, yielding the stated action.
\end{proof}

\section{Further Directions}

\begin{enumerate}
\item Classification of BV--quantized L--systems by lamphrodynamic
entropy.
\item Derived deformation theory of production rules.
\item Applications to biological growth and semantic evolution.
\end{enumerate}

\newpage

